
\subsubsection{4.1 Research Questions}

\textbf{RQ1:} Does execution test feedback achieve a statistically significantly larger pass@1 improvement than pylint/mypy at fixed B=1000 output tokens?

\textbf{RQ2:} What fraction of HumanEval baseline failures does pylint/mypy detect, and which flag categories dominate?

\textbf{RQ3:} Does the execution feedback advantage vary between HumanEval and MBPP?

\subsubsection{4.2 Datasets}

\textbf{HumanEval} \citep{Chen2021}: 164 algorithmic programming problems. Baseline pass@1 under Llama 3.1 8B Instruct with greedy decoding: 61.0\% (64/164 failures, 39\% failure rate).

\textbf{MBPP} \citep{Austin2021}: 378 function-completion problems in EvalPlus format. Baseline pass@1: 33.1\% (253/378 failures, 67\% failure rate).

\begin{table*}[htbp]
\centering
\resizebox{\dimexpr 2\columnwidth+\columnsep\relax}{!}{%
\begin{tabular}{lllll}
\toprule
Dataset & Problems & Baseline pass@1 & Failures & Problem type \\
\midrule
HumanEval & 164 & 61.0\% & 64 & Algorithmic reasoning \\
MBPP & 378 & 33.1\% & 253 & Function completion \\
\bottomrule
\end{tabular}
}
\end{table*}

\subsubsection{4.3 Evaluation Metrics}

\textbf{Primary:} $\Delta _pass$@1 per condition; McNemar's test ($\alpha =0.05)$ on paired outcomes.

\textbf{Mechanism:} Pylint total and functional (E+W) coverage fraction over HumanEval baseline failures. Bootstrap 95\% CI (10,000 samples, seed=42).

\textbf{Secondary:} Per-round pass@1 trajectory (rounds 0--3) to characterize budget saturation.

All 542 problems (164 HumanEval + 378 MBPP) were processed to completion in each condition (100\% completion rate verified in experiment logs).

